Cada \dfn{instancia} de examen atraviesa una secuencia de diez estados
que reflejan su progreso desde la llegada de la primera página hasta el
archivo definitivo del expediente. La \cref{fig:scdee-exam-state}
muestra la máquina de estados completa, incluyendo las transiciones de
retroceso permitidas para corregir errores en cualquiera de las fases.

\begin{figure}[Máquina de estados de \texttt{ExamInstance}]%
              {fig:scdee-exam-state}%
              {Ciclo de vida de una \dfn{instancia} de examen. Las
               flechas marcadas como \texttt{<<retroceso>>} representan
               las transiciones de marcha atrás permitidas al
               coordinador para corregir errores en cualquier fase del
               proceso.}
  \image{\textwidth}{}{plantuml/exam_state}
\end{figure}

Las transiciones se validan contra una tabla central
(\texttt{VALID\_TRANSITIONS}) antes de aplicarse, lo que impide saltos
arbitrarios y centraliza la lógica de cambio de estado en un único
punto. Por ejemplo, la transición a \texttt{PUBLISHED} requiere que la
instancia no tenga incidencias activas y que todos los problemas tengan
una calificación numérica asignada; cualquier intento de publicar una
instancia con incidencias devuelve un error \texttt{409 Conflict} y deja
el estado anterior intacto. Las transiciones de retroceso
—\texttt{GRADED} $\to$ \texttt{PENDING\_GRADING}, \texttt{PUBLISHED}
$\to$ \texttt{GRADED} y \texttt{FINALIZED} $\to$
\texttt{PENDING\_REVIEW}— permiten al coordinador reabrir la
corrección, retirar una publicación o atender solicitudes de revisión
tardías sin necesidad de borrar y recrear la instancia.

La fase de revisión, comprendida entre \texttt{IN\_REVIEW} y
\texttt{FINALIZED}, está gobernada por una ventana de revisión
configurada por el coordinador en el momento de la publicación. La
apertura y el cierre de la ventana se automatizan mediante Celery Beat:
al inicio de la ventana, todas las instancias en \texttt{PUBLISHED}
transitan a \texttt{IN\_REVIEW} y se notifica a los alumnos; al
cierre, las instancias sin solicitud activa transitan a
\texttt{FINALIZED}, mientras que las que tienen una solicitud abierta
quedan en \texttt{PENDING\_REVIEW} hasta que el corrector la resuelva.
El detalle del flujo de revisión —creación de solicitudes por el
alumno, resolución por el corrector y notificaciones derivadas— se
diagrama en la \cref{fig:seq-review} (\cref{AP:SEQUENCE}).

A nivel de página, cada \texttt{ExamPage} mantiene un estado más
sencillo (\texttt{PENDING\_RECOGNITION}, \texttt{RECOGNIZED},
\texttt{ORPHAN}, \texttt{DISCARDED}) independiente del estado de la
instancia a la que esté asociada. Esta separación permite reintentar
manualmente el reconocimiento de una página problemática sin afectar al
resto de la instancia, y conservar páginas huérfanas
(\texttt{ORPHAN})~\textemdash{}páginas reconocidas que no encajan en
ningún examen vigente\textemdash{} para diagnóstico posterior.
